\documentclass[11pt]{article}
% BEGIN SHARED COURSE STYLE
% Shared Math 615 style. Embedded verbatim in standalone published sources.
\RequirePackage[letterpaper,margin=1in]{geometry}
\RequirePackage{amsmath,amsthm,mathtools}
\RequirePackage{fontspec,unicode-math}
\setmainfont{texgyretermes}[Extension=.otf,UprightFont=*-regular,BoldFont=*-bold,ItalicFont=*-italic,BoldItalicFont=*-bolditalic]
\setmathfont{texgyretermes-math.otf}
\RequirePackage{microtype,tikz-cd,enumitem,needspace,xcolor,booktabs,titlesec,etoolbox}
\definecolor{teal}{HTML}{176568}
\definecolor{burgundy}{HTML}{782F40}
\definecolor{inklink}{HTML}{176568}
\definecolor{accent}{HTML}{176568}
\titleformat{\section}{\large\bfseries\color{teal}}{\thesection}{.65em}{}
\titleformat{\subsection}{\normalsize\bfseries\color{burgundy}}{\thesubsection}{.65em}{}
\titleformat{\paragraph}[runin]{\normalsize\bfseries\color{burgundy}}{}{0pt}{}
\titlespacing*{\section}{0pt}{2.2ex plus .5ex minus .2ex}{1ex plus .2ex}
\titlespacing*{\subsection}{0pt}{1.6ex plus .4ex minus .2ex}{.7ex plus .2ex}
\RequirePackage{hyperref,bookmark}
\hypersetup{colorlinks=true,linkcolor=teal,urlcolor=teal,citecolor=teal,pdfstartview={XYZ null null 1.5},bookmarksnumbered=true}
\urlstyle{same}
\setcounter{tocdepth}{2}
\setlist[enumerate]{label=\arabic*.,ref=\arabic*,leftmargin=*,itemsep=4pt,topsep=4pt}
\setlength{\parindent}{1em}
\setlength{\parskip}{3pt}
\setlength{\emergencystretch}{2em}
\raggedbottom
\AddToHook{cmd/section/before}{\Needspace{8\baselineskip}}
\AddToHook{cmd/subsection/before}{\Needspace{6\baselineskip}}
\makeatletter
\newcommand{\afterblock}{\@afterindentfalse\@afterheading}
\makeatother
\newcounter{result}[section]
\renewcommand{\theresult}{\thesection.\arabic{result}}
\newcommand{\resulttitle}[3]{#1 #2\ifstrempty{#3}{}{ (#3)}.}
\newenvironment{result}[3]{\par\addvspace{7pt}\Needspace{4\baselineskip}\refstepcounter{result}\hypertarget{#2}{}\label{#2}\noindent{\color{burgundy}\hypersetup{linkcolor=burgundy}\hyperlink{proof:#2}{\textbf{\resulttitle{#1}{\theresult}{#3}}}}\quad\ignorespaces}{\par\addvspace{4pt}\afterblock}
\newenvironment{entry}[3]{\par\addvspace{7pt}\Needspace{3\baselineskip}\refstepcounter{result}\hypertarget{#2}{}\label{#2}\noindent{\color{burgundy}\textbf{\resulttitle{#1}{\theresult}{#3}}}\quad\ignorespaces}{\par\addvspace{4pt}\afterblock}
\newenvironment{demonstration}[1]{\par\noindent\hypertarget{proof:#1}{}{\color{burgundy}\hypersetup{linkcolor=burgundy}\hyperlink{#1}{\textbf{Proof.}}}\quad\ignorespaces}{\unskip\nobreak\hfill\qedsymbol\par\addvspace{5pt}\afterblock}
\newenvironment{citedresult}[4]{\par\addvspace{7pt}\Needspace{4\baselineskip}\refstepcounter{result}\hypertarget{#2}{}\label{#2}\noindent{\color{burgundy}\hypersetup{urlcolor=burgundy}\href{#4}{\textbf{\resulttitle{#1}{\theresult}{#3}}}}\quad\ignorespaces}{\par\addvspace{4pt}\afterblock}
\newcommand{\usedby}[1]{\par\noindent{\small\textbf{Used in:} #1.}\par\afterblock}
\newcommand{\coursebold}[1]{{\color{burgundy}\textbf{#1}}}
\newcommand{\coursetitle}[3]{\begin{center}{\large\bfseries\color{teal}#1}\\[4pt]{\LARGE\bfseries\color{teal}#2}\\[6pt]Math 615: Algebraic Topology I\\Fall 2026\quad\textbar\quad #3\end{center}}
\newcommand{\coursecontents}{{\small\setlength{\parskip}{0pt}\tableofcontents}}
\newcommand{\homeworkstyle}{\titleformat{\section}{\large\bfseries\color{teal}}{Exercise \thesection.}{.65em}{}}
% END SHARED COURSE STYLE
\newcommand{\Z}{\mathbb Z}
\newcommand{\R}{\mathbb R}
\newcommand{\id}{\mathrm{id}}
\newcommand{\sSet}{\mathcal P(\Delta)}
\newcommand{\sAb}{\mathcal P(\Delta;\mathrm{Ab})}
\newcommand{\Top}{\mathrm{Top}}
\newcommand{\Sing}{\operatorname{Sing}}
\newcommand{\Mor}{\operatorname{Mor}}
\newcommand{\im}{\operatorname{im}}
\newcommand{\C}{C_\bullet}
\homeworkstyle
\hypersetup{pdftitle={Homework 4: Simplicial Sets and Singular Chains},pdfauthor={Math 615 Fall 2026},pdfsubject={Revision 3}}
\begin{document}
\coursetitle{Homework 4}{Simplicial Sets and Singular Chains}{September 21 and 23}
Use the constructions and results of Lecture 4. Give proofs and justify
all asserted naturality statements. Unless a quotient is explicitly specified,
retain all simplices, including degenerate ones. The definitions below fix the
notation; the exercises begin on page~\pageref{ex:standard}.

\section*{Definitions and Notation}
\coursebold{Functor Categories and Presheaves.}
The category $\operatorname{Fun}(\mathcal C,\mathcal D)$ has functors
$\mathcal C\to\mathcal D$ as objects and natural transformations as morphisms.
A natural transformation $\theta:F\to G$ consists of morphisms
$\theta_c:F(c)\to G(c)$ satisfying
$G(f)\theta_c=\theta_{c'}F(f)$ for every $f:c\to c'$.
The opposite category $\mathcal C^{\mathrm{op}}$ reverses all arrows of
$\mathcal C$. Write
\[
\mathcal P(\mathcal C;\mathcal D)=\operatorname{Fun}(\mathcal C^{\mathrm{op}},\mathcal D),
\qquad
\mathcal P(\mathcal C)=\mathcal P(\mathcal C;\mathrm{Set}).
\]
Objects of $\mathcal P(\mathcal C)$ are called \emph{presheaves of sets}.
Here $\mathrm{Set}$, $\mathrm{Ab}$, and $\Top$ denote the categories of sets,
abelian groups, and topological spaces, with functions, homomorphisms, and
continuous maps respectively. The notation $\Mor_{\mathcal C}(c,c')$ means
the set of morphisms in $\mathcal C$.

\coursebold{The Simplex Category and Simplicial Objects.}
The category $\Delta$ has objects $[n]=\{0<\cdots<n\}$, $n\geq0$,
and weakly increasing maps as morphisms. A simplicial set is a presheaf
$K\in\mathcal P(\Delta)$; write $K_n=K([n])$ and
$\alpha^*=K(\alpha):K_n\to K_m$ for $\alpha:[m]\to[n]$.
An element of $K_n$ is an \emph{$n$-simplex}; vertices, edges, and triangles
are simplices of degrees zero, one, and two. A simplicial map $f:K\to L$
is a natural transformation, so $f_m\alpha^*=\alpha^*f_n$.
A simplicial abelian group is an object of
$\mathcal P(\Delta;\mathrm{Ab})$: its restriction maps are homomorphisms.

\coursebold{Faces, Degeneracies, and the Standard Simplex.}
For $n\geq1$ and $0\leq i\leq n$, let
$\delta^i:[n-1]\hookrightarrow[n]$ be the injection omitting $i$.
For $n\geq0$ and $0\leq i\leq n$, let
$\sigma^i:[n+1]\twoheadrightarrow[n]$ be the surjection identifying
$i$ and $i+1$. Thus
\[
\delta^i(k)=\begin{cases}k&k<i,\\k+1&k\geq i,\end{cases}
\qquad
\sigma^i(k)=\begin{cases}k&k\leq i,\\k-1&k>i.\end{cases}
\]
The induced maps $d_i=(\delta^i)^*:K_n\to K_{n-1}$ and
$s_i=(\sigma^i)^*:K_n\to K_{n+1}$ are the \emph{face} and
\emph{degeneracy} maps. For $n>0$, a simplex of $K_n$ is
\emph{degenerate} if it lies in some $s_i(K_{n-1})$, and
\emph{nondegenerate} otherwise; every vertex is nondegenerate.
The standard simplicial $n$-simplex is
\[
\Delta^n_m=\Mor_\Delta([m],[n]),\qquad
\beta^*(\alpha)=\alpha\circ\beta.
\]
A map $\alpha:[m]\to[n]$ also induces a map
$\Delta^m\to\Delta^n$ by postcomposition.


\coursebold{Barycentric Coordinates.}
The standard topological $n$-simplex is
\[
|\Delta^n|=\{(t_0,\ldots,t_n)\in\R^{n+1}:
 t_i\geq0,\ t_0+\cdots+t_n=1\},
\]
with the subspace topology from $\R^{n+1}$.
Its vertices are the standard basis vectors $e_0,\ldots,e_n$.
Each point has the unique expression $t_0e_0+\cdots+t_ne_n$;
the coefficients $t_i$ are its \emph{barycentric coordinates}.
For example,
$(1/2,1/3,1/6)$ is a point of $|\Delta^2|$, while $(1/2,1/2,0)$
lies on the edge opposite $e_2$. The barycenter has every coordinate
$1/(n+1)$.

For a weakly increasing map $\alpha:[m]\to[n]$, the associated affine map
$\alpha_*:|\Delta^m|\to|\Delta^n|$ sends $e_i$ to $e_{\alpha(i)}$:
\[
(\alpha_*t)_j=\sum_{\alpha(i)=j}t_i
\quad(\text{an empty sum is }0).
\]
In particular, $(\delta^i)_*$ inserts a zero coordinate and
$(\sigma^i)_*$ adds adjacent coordinates $t_i+t_{i+1}$.
A formula ``in barycentric coordinates'' is a formula in these $t_i$.

\coursebold{Realization and Singular Simplices.}
The geometric realization of $K$ is the quotient space
\[
|K|=\left(\coprod_{n\geq0}K_n\times|\Delta^n|\right)/\sim,
\qquad
(\alpha^*x,t)\sim(x,\alpha_*t),
\]
where each $K_n$ is discrete and the displayed pairs generate the
equivalence relation. Write $[x,t]$ for the class of $(x,t)$.
The characteristic map of $x\in K_n$ is
$\chi_x:|\Delta^n|\to|K|$, $t\mapsto[x,t]$.
A \emph{singular $n$-simplex} in $X$ is a continuous map
$\tau:|\Delta^n|\to X$. These form the simplicial set $\Sing X$,
with $\alpha^*\tau=\tau\circ\alpha_*$.
The realization--singular adjunction is the natural bijection
\[
\Mor_\Top(|K|,X)\cong\Mor_{\mathcal P(\Delta)}(K,\Sing X).
\]
Its \emph{unit} $\eta_K:K\to\Sing|K|$ sends $x$ to $\chi_x$.

\coursebold{Subsets, Boundaries, and Quotients.}
A simplicial subset $L\subseteq K$ consists of subsets $L_n\subseteq K_n$
preserved by every restriction map. The boundary
$\partial\Delta^n\subseteq\Delta^n$ consists in degree $m$ of the
nonsurjective maps $[m]\to[n]$. When each $L_n$ is nonempty, the
quotient $K/L$ identifies all elements of $L_n$ to one element in each
degree; its operators are induced from those of $K$.
The notation $B\amalg_A C$ denotes a \emph{pushout}: an object receiving
maps from $B$ and $C$ agreeing on $A$, universal among such objects.
Coproducts $K\amalg L$, products $K\times L$, and pushouts of simplicial
sets are computed degree by degree in sets.


\section{Standard Simplices and Simplicial Operators}\label{ex:standard}
Represent an $m$-simplex of $\Delta^n$ by a weakly increasing sequence
$(a_0,\ldots,a_m)$ with $0\leq a_i\leq n$.
\begin{enumerate}
\item Describe every face and degeneracy map on these sequences. Verify
$d_i d_j=d_{j-1}d_i$ for $i<j$, and all three cases of the identity for
$d_i s_j$, by tracking the entries deleted or repeated.
\item Prove that a simplex is nondegenerate precisely when its entries
are strictly increasing. Count the $m$-simplices and the nondegenerate
$m$-simplices of $\Delta^n$. List the edges of $\Delta^2$, distinguishing
its degenerate and nondegenerate edges.
\item For $x\in K_n$, define a map $\bar x:\Delta^n\to K$ by
$\bar x_m(\alpha)=\alpha^*x$. Prove that evaluation at $\id_{[n]}$
gives a natural bijection $\Mor_{\sSet}(\Delta^n,K)\cong K_n$.
Under this bijection, identify precomposition by the maps of standard
simplices induced by $\delta^i$ and $\sigma^i$.
\end{enumerate}

\section{Singular Simplices and Degeneracy}\label{ex:singular}
The homeomorphism $[0,1]\to|\Delta^1|$ is $u\mapsto(1-u,u)$.
For a space $X$, recall that $(\Sing X)_n=\Mor_\Top(|\Delta^n|,X)$.
\begin{enumerate}
\item For a singular triangle $\tau:|\Delta^2|\to X$, write formulas
for $d_0\tau,d_1\tau,d_2\tau$ as paths parametrized by $u\in[0,1]$.
Identify the initial and terminal vertex of each path.
\item For a path $\gamma:[0,1]\to X$, write $s_0\gamma$ and
$s_1\gamma$ in barycentric coordinates, and compute all three faces
of each. Specify which resulting paths are constant.
\item Prove that a singular edge is degenerate if and only if it is
constant. Give an explicit nonconstant path in $\R$ whose two endpoints
are equal. Explain why it is nevertheless nondegenerate.
\item Describe the image in degrees zero and one of the unit
$\eta_{\Delta^1}:\Delta^1\to\Sing|\Delta^1|$. Give a singular vertex
and a singular edge outside the respective images. Deduce that this
unit is not an isomorphism of simplicial sets.
\end{enumerate}


\section{Realization, Quotients, and Loops}\label{ex:realization}
Let $S=\Delta^1/\partial\Delta^1$ be the degreewise quotient collapsing
the simplicial subset generated by the two vertices to a single vertex.
Write $v$ for its vertex and $e$ for its nondegenerate edge.
\begin{enumerate}
\item Describe $S_m$ in terms of weakly increasing sequences of zeros
and ones. Prove that $v$ and $e$ are its only nondegenerate simplices.
List all simplices in degrees one and two, including degeneracies.
\item Use the realization--singular adjunction to prove that realization
preserves pushouts. Apply this to
$S=\Delta^1\amalg_{\partial\Delta^1}\Delta^0$ to identify $|S|$
with $[0,1]/(0\sim1)$, including its topology. Explain why the
simplicial quotient introduces no extra geometric cells.
\item Prove that a simplicial map $S\to K$ is equivalent to the choice
of $x\in K_0$ and $a\in K_1$ with $d_0a=d_1a=x$.
Apply this to $K=\Sing X$ and identify such maps with pairs consisting
of a point $x\in X$ and a continuous loop $\gamma:[0,1]\to X$
based at $x$. These are actual loops, not homotopy classes.
\item For this correspondence, write explicitly the associated continuous
map $|S|\to X$. Check that it agrees with the adjunction formula
the formula $F([z,t])=g_n(z)(t)$ for a simplicial map $g:S\to\Sing X$.
\end{enumerate}

\section{Free Simplicial Abelian Groups}\label{ex:free}
For a set $T$, the \emph{free abelian group} $\Z[T]$ consists of finite
formal sums $\sum_{t\in T}n_t t$ with $n_t\in\Z$.
A function $T\to U(B)$ to the underlying set of an abelian group extends
uniquely to a homomorphism $\Z[T]\to B$; this is the free--forgetful
adjunction. For $K\in\sSet$, put $\Z[K]_m=\Z[K_m]$, with restriction
maps obtained by extending those of $K$ linearly. A generator associated to
$x\in K_m$ is denoted by the same symbol $x$. Let $A\in\sAb$, and let
$U(A)$ be its underlying simplicial set, obtained by forgetting the group structure in each degree.
For an adjunction $F\dashv U$, its unit $\eta:\id\to UF$ and counit
$\epsilon:FU\to\id$ correspond to identity morphisms under the
adjunction bijections. Their triangle identities are
$\epsilon_{F(K)}\circ F(\eta_K)=\id_{F(K)}$ and
$U(\epsilon_A)\circ\eta_{U(A)}=\id_{U(A)}$.
\begin{enumerate}
\item Combine the free--forgetful adjunction with
Exercise~\ref{ex:standard} to construct a natural group isomorphism
\[
\Mor_{\sAb}(\Z[\Delta^n],A)\cong A_n.
\]
Describe the homomorphism associated to $a\in A_n$ on every generator
$\alpha$ in every degree. The group structure on the left
is pointwise addition of simplicial homomorphisms.
\item For $S$ from Exercise~\ref{ex:realization}, identify
$\Mor_{\sAb}(\Z[S],A)$ with the subgroup
\[
\{(x,a)\in A_0\times A_1:d_0a=d_1a=x\}.
\]
Describe directly the values of the associated homomorphism on
$ s_0v$, $ s_0e$, and
$ s_1e$.
\item Write the unit $K\to U\Z[K]$ and counit $\Z[U(A)]\to A$.
Verify the two triangle identities on generators. In $\Z[U(A)]_n$,
exhibit a nonzero element sent to zero by the counit; explain why
$0$ is not the zero element of this free abelian group.
\item Prove naturally that $\Z[K\amalg L]\cong\Z[K]\oplus\Z[L]$.
For $K=L=\Delta^0$, show that the canonical map induced by the two
projections,
$\Z[K\times L]\to\Z[K]\times\Z[L]$, is not an isomorphism.
\end{enumerate}


\section{Alternating Faces and Degenerate Chains}\label{ex:chains}
For a simplicial abelian group $A$, let $C_n(A)=A_n$ and
$\partial_n=\sum_{i=0}^n(-1)^i d_i$ for $n>0$.
For a simplicial set $K$, abbreviate $C_\bullet(\Z[K])$ to
$C_\bullet(K;\Z)$. Set $C_n=0$ for $n<0$ and $\partial_0=0$.
A chain complex has $\partial_{n-1}\partial_n=0$; a chain map commutes
with the differentials. A subcomplex $D_\bullet\subseteq C_\bullet$
means $\partial_n(D_n)\subseteq D_{n-1}$, so the quotient groups carry
an induced differential. The \emph{unnormalized} complex retains all
simplices; the quotient by degenerate chains is called the
\emph{normalized quotient complex}.
\begin{enumerate}
\item Expand $\partial_2\partial_3$ and pair its twelve terms using the
face identities. Explain how the same pairing proves $\partial^2=0$
in every degree. Show that a simplicial homomorphism induces a chain map.
\item Put $D_0(A)=0$ and
$D_n(A)=\sum_{j=0}^{n-1}\im(s_j:A_{n-1}\to A_n)$ for $n>0$.
For $a\in A_{n-1}$, expand $\partial_n(s_ja)$ using the mixed
face--degeneracy identities. Show that the two terms not visibly
in a degeneracy image cancel. Deduce that $D_\bullet(A)$ is a
subcomplex, treating $n=1$ explicitly.
\item For $A=\Z[K]$, prove that $D_n(A)$ is freely generated by the
degenerate $n$-simplices. Hence identify $C_n(A)/D_n(A)$ with the free
abelian group on the nondegenerate $n$-simplices, and describe its
induced differential. You need not prove that this quotient preserves
homology.
\item Compute the \emph{unnormalized} complex $C_\bullet(\Delta^0;\Z)$
and its homology in all degrees. For its vertex $v$, explain explicitly
why $ s_0v$ is a nonzero cycle but represents zero in
homology. Compare with the quotient complex from the preceding part.
\end{enumerate}

\section{Components and a First Circle Calculation}\label{ex:homology}
Write $H_n(K;\Z)=\ker(\partial_n)/\im(\partial_{n+1})$ for the
homology of $C_\bullet(K;\Z)$: elements of the kernel are \emph{cycles}
and elements of the image are \emph{boundaries}.
Define $H_n^{\mathrm{sing}}(X;\Z)=H_n(\Sing X;\Z)$.
A path from $x$ to $y$ is a continuous map $\gamma:[0,1]\to X$
with $\gamma(0)=x$, $\gamma(1)=y$; it is a loop based at $x$ if $y=x$.
Path components are equivalence classes under the existence of a path;
connected components are maximal connected subsets.
Let $\pi_0(K)$ denote the quotient of $K_0$ by the equivalence relation
generated by $d_1a\sim d_0a$ for every $a\in K_1$.
\begin{enumerate}
\item Prove that the map taking each vertex to its equivalence class
induces a natural isomorphism
\[
H_0(K;\Z)\cong\Z[\pi_0(K)].
\]
Show that the kernel of $\Z[K_0]\to\Z[\pi_0(K)]$ is generated by
$ d_0a- d_1a$.
\item Deduce naturally that
$H_0^{\mathrm{sing}}(X;\Z)\cong\Z[\pi_0^{\mathrm{path}}(X)]$,
where $\pi_0^{\mathrm{path}}(X)$ is the set of path components.
Explain why connected components cannot simply replace path components
in this argument. No example of a connected, non-path-connected space
is required.
\item If $X$ is discrete, identify $\Sing X$ and compute
$H_n^{\mathrm{sing}}(X;\Z)$ for every $n\geq0$, including $X=\varnothing$.
Use the actual unnormalized differential.
\item For $S=\Delta^1/\partial\Delta^1$, use the lists from
Exercise~\ref{ex:realization} to compute $C_0(S;\Z)$, $C_1(S;\Z)$,
$C_2(S;\Z)$ and both differentials. Deduce $H_0(S;\Z)$ and
$H_1(S;\Z)$, specifying a generator of $H_1$. Work directly in the
unnormalized complex; do not assume a comparison theorem with the
singular homology of $|S|$.
\end{enumerate}
\end{document}
